\section{System}
\label{sec:system}

\begin{figure*}[t]
  \centering
  \resizebox{\textwidth}{!}{%
  \begin{tikzpicture}[
      node distance=3mm and 5mm,
      box/.style={draw, rounded corners=2pt, align=center, font=\scriptsize,
                  minimum height=8mm, inner sep=3pt, fill=#1},
      box/.default=white,
      lab/.style={font=\scriptsize\itshape, text=black!60},
      arr/.style={-{Latex[length=1.6mm]}, thick, black!70}]
    % ---------------- offline
    \node[box=gray!10, minimum width=18mm] (corpus) {domain corpus\\(1.17M documents)};
    \node[box=blue!8, right=of corpus, yshift=7mm, minimum width=26mm] (chunk)
      {200-word passages\\(2.41M, MaxP)};
    \node[box=blue!8, right=of corpus, yshift=-12mm, minimum width=26mm] (meta)
      {Qwen3-1.7B metadata:\\period, events, entities};
    \node[box=blue!15, right=of chunk, yshift=7mm, minimum width=24mm] (bm25) {BM25};
    \node[box=blue!15, below=1.5mm of bm25, minimum width=24mm] (dense)
      {Qwen3-Embedding-0.6B};
    \node[box=blue!15, below=1.5mm of dense, minimum width=24mm] (rir) {ReasonIR-8B};
    \node[box=blue!15, below=1.5mm of rir, minimum width=24mm] (mindex) {metadata index};
    \draw[arr] (corpus) -- (chunk); \draw[arr] (corpus) -- (meta);
    \draw[arr] (chunk.east) -- (bm25.west); \draw[arr] (chunk.east) -- (dense.west);
    \draw[arr] (chunk.east) -- (rir.west); \draw[arr] (meta.east) -- (mindex.west);
    \node[lab, above=0.5mm of chunk] {offline, per domain};
    % ---------------- online
    \node[box=gray!10, right=14mm of bm25, minimum width=16mm] (q) {question\\or step};
    \node[box=orange!15, below=5mm of q, minimum width=30mm] (llm)
      {Qwen3-8B (+ top-5 BM25):\\pseudo answer, sub-queries,\\neighbouring events,
       reasoning};
    \node[lab, above=0.5mm of q] {online};
    \node[box=orange!8, right=8mm of q, yshift=-6mm, minimum width=30mm] (chan)
      {first-stage channels\\(BM25, dense, ReasonIR\\with the expanded queries)};
    \node[box=green!12, right=of chan, minimum width=22mm] (pool)
      {pool: weighted RRF\\of channels chosen\\on train};
    \node[box=green!20, right=of pool, minimum width=26mm] (rr)
      {Qwen3-Reranker\\scores the top-100};
    \node[box=green!28, right=of rr, minimum width=12mm] (fin) {RRF};
    \node[box=gray!10, right=of fin, minimum width=10mm] (out) {top-10};
    \draw[arr] (q) -- (llm);
    \draw[arr] (q.east) -- (chan.north west);
    \draw[arr] (llm.east) -- (chan.south west);
    \draw[arr] (chan) -- (pool); \draw[arr] (pool) -- (rr); \draw[arr] (rr) -- (fin);
    \draw[arr] (pool.south) -- ++(0,-5mm) -| (fin.south);
    \draw[arr] (fin) -- (out);
    \draw[arr, dashed] (bm25.east) -- (q.west);
    \draw[arr, dashed] (mindex.east) -| (chan.south);
  \end{tikzpicture}}
  \caption{\sys. Offline, every document is split into passages, indexed with BM25 and two
  dense encoders, and summarised by a small LLM into a metadata record. Online, Qwen3-8B
  expands the question (given the top BM25 results); the expanded queries run through the
  indexes, a train-selected subset of these channels is fused into a pool, and a decoder
  reranker scores the pool's top-100, which is fused with the pool into the final top-10.}
  \label{fig:overview}
\end{figure*}

Figure~\ref{fig:overview} summarises \sys. It has three parts: indexes built once per domain
(\S\ref{sec:indexes}), first-stage channels that combine those indexes with LLM-expanded
queries (\S\ref{sec:channels}), and a train-selected pool reranked by a decoder LLM
(\S\ref{sec:pool}). Every component runs per domain and is cached, so all variants in
Table~\ref{tab:full} are evaluated on identical candidates.

\subsection{Indexes}
\label{sec:indexes}
\paragraph{Passages.} TEMPO documents range from a sentence to multi-megabyte dumps. We split
every document into 200-word passages with a 40-word overlap (at most 64 per document; 2.41M
passages in total) and score a document by its best passage (MaxP, \citealp{dai2019maxp}).
Passages after the first are prefixed with the document's title, so that they keep its topic.

\paragraph{BM25.} Lucene BM25 with the official parameters ($k_1{=}0.9$, $b{=}0.4$) over
passages, built with \texttt{bm25s} and a streaming index builder that reproduces its scores
exactly and fits the largest domain (\emph{history}, 200k documents) in memory.

\paragraph{Dense.} All passages are embedded with Qwen3-Embedding-0.6B
\citep{zhang2025qwen3emb} and, separately, with ReasonIR-8B \citep{shao2025reasonir}, a
Llama-3.1-8B encoder trained for reasoning-intensive retrieval. Search is exact (inner product
on GPU) followed by MaxP; for Qwen3-Embedding we also keep document vectors (the mean of a
document's passage vectors).

\paragraph{LLM metadata.} A small LLM (Qwen3-1.7B, \citealp{yang2025qwen3}) reads the opening of
every document (220 words, at most 1{,}600 characters) and returns, as JSON, the period the text
is \emph{about}---including implicit periods such as ``after the 2008 crisis''---up to five
events, six entities, six keywords and a subtopic. The events, entities, keywords and title of
each document form a short text that is indexed with BM25 at document level (the
\emph{metadata index}). The same prompt is applied to every question, so questions and
documents share one schema.

\subsection{Query Expansion and Channels}
\label{sec:channels}
All expansions come from Qwen3-8B \citep{yang2025qwen3} in non-thinking mode with greedy,
JSON-schema-guided decoding. Following LLM-based pseudo-relevance feedback, the prompts can
include the first 60 words of the top-5 BM25 documents, so that the model reuses the
vocabulary, names and dates the collection actually uses. Each expansion becomes a
\emph{channel}, a ranked list of 100 documents:

\begin{itemize}
  \setlength\itemsep{1pt}
  \item \textbf{query2doc}: BM25 with the question plus a 60--90-word hypothetical answer and
        8--15 keywords \citep{wang2023query2doc}, with or without BM25 feedback. A single
        demonstration, built from one train question of the same domain (its released
        guidance and its first gold document), shows the model the expected answer style.
  \item \textbf{Sub-queries}: the same prompt also returns three sub-queries and a
        chronological event chain; one BM25 run per sub-query and per event is fused with RRF, the
        question weighted twice \citep{rackauckas2024ragfusion}.
  \item \textbf{Neighbouring events}: relevant evidence often describes events \emph{around}
        the one a question names. A second prompt (with BM25 feedback) asks for 4--6 such
        events---causes, earlier, parallel and later events, each with a year---four
        period-anchored sub-queries and ten expansion terms. BM25 and Qwen3-Embedding runs for
        the question (weight 2), for every sub-query and for every event are fused with RRF.
  \item \textbf{ReasonIR}: the full question (for a step: question and step), and a second
        variant that appends Qwen3-8B's free-text reasoning about the question, as done with
        GPT-4o on TEMPO \citep{abdallah2026tempo}.
  \item \textbf{Dense} (Qwen3-Embedding, passages and document vectors) and the
        \textbf{metadata index}, queried with the question and its own small-LLM metadata.
\end{itemize}

\subsection{Pool and Reranking}
\label{sec:pool}
\paragraph{Pool.} Channels are combined with weighted reciprocal rank fusion
\citep{cormack2009rrf}, $s(d){=}\sum_c w_c/(60{+}\mathrm{rank}_c(d))$. Which channels enter, and
with which weight, is decided on the \emph{train} split by greedy forward selection: starting
from the best single channel, the (channel, $w\in\{1,\tfrac12\}$) pair that most improves train
macro nDCG@10 is added until no addition helps. Selection is run separately for 1a and 1b, over
all channels of Table~\ref{tab:full}.

\paragraph{Reranking.} The pool's top-100 is reranked pointwise and the reranked list is fused
with the pool by equal-weight RRF. Each candidate is shown as ``[first--last year] title $|$
best passage'', the best passage being the document's highest-scoring BM25 passage for the
question. We compare encoder \emph{cross-encoders}---MiniLM-L12 trained on MS~MARCO
\citep{wang2020minilm,nogueira2019passage} and BGE-reranker-v2-m3 \citep{chen2024bgem3}---with
decoder LLM rerankers, \emph{cross-decoders}: Qwen3-Reranker 0.6B and 4B \citep{zhang2025qwen3emb},
which read an instruction, the question and the document and score $P(\textit{yes})$ at the
answer position \citep[cf.][]{nogueira2020monot5}. Their instruction states RETECO's notion of
relevance (``evidence that helps answer it, including information about the relevant events and
time period''). The reranker is the one with the best train score for the sub-track;
reranking the BM25 $\oplus$ dense candidates instead of the pool's scored lower on train for the
selected rerankers.

\paragraph{Sub-track 1b.} A step is retrieved with ``title + step instruction'' (ReasonIR: full
question + step). The expansions, metadata and reasoning of the parent question are shared by
all of its steps.

\paragraph{Time-matching ablations.} To test whether explicit time helps, we also build
components that are \emph{not} part of the final system. A rule-based extractor normalises
years, ranges and dates in every document (in the final system it only supplies the year span
shown to the reranker); from it we build a year-filtered BM25, a year prior
$s'(d){=}(1{-}\lambda)\hat s(d){+}\lambda T(q,d)$ ($T$: fraction of the question's years the
document mentions, $\pm1$), and a document--entity--year graph searched with personalised
PageRank; with the LLM metadata, a period filter for BM25 and dense retrieval (overlap with the
question's period, $\pm2$ years, undated documents kept) and a period prior. $\lambda$ is chosen
on train. We also run RM3 pseudo-relevance feedback \citep{lavrenko2001relevance,abduljaleel2004umass}.
